\paragraph{Comparison of ground-truth forms.}
Taken together, these works instantiate three complementary forms of explanation
ground truth. The benchmarks of \citet{baric2021benchmarking},
\citet{almaghrabi2024multidimensional}, and
\citet{delibasoglu2025crossscale} derive structural or binary relevance from
known variables and lags. DCIts compares learned local transition coefficients
with the exact signed coefficients of the generator. SHAPformer derives
value-level ground truth by applying a conditional Shapley-value procedure
directly to the generating function. The corresponding evaluations are,
respectively, localization or ranking against structural masks, signed
coefficient recovery, and comparison of Shapley attribution
sign, magnitude, and forecast-horizon dependence.

\paragraph{Incremental contribution of each benchmark.}
Table~\ref{tab:incremental-gt-contributions} orders the works chronologically
and reports only the capability introduced beyond the rows above it. The label
\emph{explainable by design} identifies a model that exposes an explanation as
part of its own computation. A direct mechanism-recovery error additionally
requires this exposed object to have the same mathematical meaning and indexing
as the generator mechanism. Consequently, intrinsic attention and saliency
scores can be evaluated against support locations or relevant regions, while an
explicit transition tensor such as DCIts can be compared coefficient by
coefficient. A post-hoc attribution can also be compared with process-level
ground truth; that result measures the combined ability of the forecaster to
learn the process and of the explainer to describe the forecaster, unless a
separate model-level reference explanation is available.

\begingroup
\footnotesize
\setlength{\LTpre}{0.5\baselineskip}
\setlength{\LTpost}{0.5\baselineskip}
\renewcommand{\arraystretch}{1.08}
\begin{longtable}{@{}
>{\raggedright\arraybackslash}p{0.18\textwidth}
>{\raggedright\arraybackslash}p{0.35\textwidth}
>{\raggedright\arraybackslash}p{0.39\textwidth}@{}}
\caption{Incremental contributions to ground-truth generation and explanation
evaluation. Each row states what is new relative to the papers listed above it.}
\label{tab:incremental-gt-contributions}\\
\toprule
Paper and ground-truth family &
New ground-truth-generation capability &
Explanation source and comparison with ground truth \\
\midrule
\endfirsthead

\multicolumn{3}{@{}l}{\small\emph{Table~\thetable\ continued}}\\
\toprule
Paper and ground-truth family &
New ground-truth-generation capability &
Explanation source and comparison with ground truth \\
\midrule
\endhead

\midrule
\multicolumn{3}{r@{}}{\small\emph{Continued on the next page}}\\
\endfoot

\bottomrule
\endlastfoot

\citet{baric2021benchmarking}

\emph{Structural binary support} &
Establishes the baseline family: ten transparent multivariate generators define
a binary target--source--lag mask from the terms that occur in each equation.
The suite varies autoregression, cross-series dependence, nonlinearity, and
switching behavior. &
The evaluated outputs are architecture-bound: seq2graph, IMV-LSTM, and DA-RNN
expose intrinsic attention; TCDF returns detected directed relations and delays.
Mean/standard-deviation heatmaps and TCDF retrieval frequencies are checked
against the known support. This introduces source and lag localization with
run-to-run stability. The reported explanation assessment consists of
model-specific visual alignment, attention statistics, and retrieval frequency;
nonnegative attention supports location and ranking, while signed coefficients
lie outside this explanation space. \\
\addlinespace

\citet{almaghrabi2024multidimensional}

\emph{Structural binary support, as in Bari\'c et al.} &
Keeps Bari\'c et al.'s binary variable--lag family and adds a forecast-horizon
axis for sequence-to-sequence prediction. The relevant locations shift by one
time step for each horizon, and nonlinear multiplicative terms test whether the
method identifies interacting inputs. &
MDA is explainable by design: its dynamic attention weights participate in
filtering the input for each prediction step. The variable and temporal scores
are multiplied (paper Eq.~12), and the resulting maps are visually checked
against known source--lag locations for Toy1 at $h=1,2,4$ and Toy2 at
$h=2,20$ (Figs.~7--8). This adds horizon-specific inspection to Bari\'c et
al.'s support-localization approach; no numerical ground-truth explanation
error is reported. MAE/RMSE and architecture ablations assess forecasting;
explanation agreement remains map-based. \\
\addlinespace

\citet{delibasoglu2025crossscale}

\emph{Structural saliency regions, in the same binary-support family as Bari\'c
et al. and MDA} &
Adds controlled multi-scale temporal regions---recent, distant, scattered, and
mixed---together with predefined relevant feature sets across SYN1--SYN8. The
ground truth therefore tests the spatial resolution and long-range localization
of a saliency method rather than coefficient magnitude. &
CrossScaleNet is explainable by design through intrinsic cross-scale attention.
As in the preceding works, predicted saliency maps are visually compared with
known support, now at the level of multi-scale regions (Figs.~8--9). Feature
Ablation and Integrated Gradients provide post-hoc feature-importance checks
on synthetic data (Fig.~5). Sufficiency and comprehensiveness add numerical,
ground-truth-free faithfulness tests on real data (Table~4). Synthetic
mask agreement remains a visual assessment. \\
\addlinespace

\citet{hertel2026shapformer}

\emph{Signed local contribution} &
Introduces a different family: conditional Shapley values are computed directly
for the vector-valued load generator. Dependency-aware resampling supplies
absent feature groups, producing signed, sample-specific contributions for each
feature group and forecast step. This adds contribution magnitude, direction,
feature dependence, interactions, and horizon variation to the earlier support
masks. &
SHAPformer is explainable by design through masked training that yields
sampling-free SHAP-style values; the experiments also use post-hoc SHAP for a
standard Transformer. Global absolute-importance bars, dependence and
interaction plots, and local contribution curves are compared with generator
SHAP. This is the first row to compare signed contribution values instead of
only support. Agreement is reported graphically across the global, dependence,
interaction, and local views. Because the reference explains the generator, the
comparison includes both process learning and explanation recovery. \\
\addlinespace

\citet{baric2026dcits}

\emph{Signed coefficient and realized local contribution} &
Reuses and extends Bari\'c et al.'s generators, then promotes the reference from
a binary occurrence mask to exact signed coefficients, biases, regimes, and
polynomial orders. Multiplying a coefficient by the observed source--lag value
also yields a sample-specific numerical contribution. &
DCIts is explainable by design and exposes the transition tensor that directly
forms its forecast. It is therefore the first row permitting a clean,
coefficient-by-coefficient mechanism-recovery comparison in the generator's own
parameter space. The paper compares selected learned and exact coefficients
with lag-resolved heatmaps and across-run mean $\pm$ standard deviation;
$|\mu|>1.95\sigma$ filters unstable entries. VAR matrices and polynomial
branches test sign, magnitude, lag, and order. Test MSE and its across-run
variation assess forecasting separately. The coefficient comparisons are
reported individually rather than combined into one full-tensor explanation
score. \\

\end{longtable}
\endgroup

